
\documentclass[a4paper,twocolumn,landscape]{article}
\usepackage{geometry}
\usepackage{concmath}
\usepackage[OT1]{fontenc}
\usepackage{graphicx}
\usepackage{amsfonts}
\usepackage{amsmath}
\usepackage{calrsfs}
\usepackage{color}
\usepackage[dvipsnames]{xcolor}
\usepackage[many]{tcolorbox}
\usepackage{setspace}               % for LINE SPACING
\usepackage{multicol}  
\geometry{margin=0.5in}
\definecolor{green}{RGB}{52, 235, 168}
\definecolor{violet}{RGB}{165, 83, 237}
\definecolor{lemon}{RGB}{170, 237, 83}
\definecolor{red}{RGB}{240, 125, 84}
\definecolor{blue}{RGB}{105, 255, 242}

\begin{document}
création à partir de rien : création ex nihilo\\
c'est un mot très chargé, on peut parler de créativité\\

création :\\
- radicalement neuf\\
- pas seulement cantonnée au champs de l'art \\

création scientifique\\
certaines créations échappent à son créateur : Frankenstein, IA, bombe atomique\\
idée très ancienne \(\rightarrow\) arcane, mystère\\

Qui crée ? \\
Il y a quelque chose entre le corporel et le spirituel \\

- Kant\\
l'oeuvre d'art n'est pas réductible à un concept  mais pas à une forme pure, c'est un jeu entre la forme et l'idée \\

l'art imite ou crée ? \\
Le rapport entre l'art, l'illusion et la vérité\\

l'art de la dialectique : l'art du dialogue \\

Ion dit que la source du savoir est la vérité pas la poésie \\
Dans Ion, Ion ne comprend pas pourquoi il ne connaît pas d'autres poètes que Homère, ni pq Homère est si bien. Socrate lui montre
qu'il n'est pas détenteur d'une tekne, mais qu'il est mu par une volonté divine 
\(\Leftrightarrow\) pierre de magnésie : rapproche les spectateurs de dieu. 
La raison de l'art n'est pas dans l'artiste, mais dans l'esprit du dieu. 

La République : cité idéale \\

L'art est dangereux, éloigne de la vérité, c'est un simulacre mais ça a de grands effets\\
Problème : agi sur l'âme humaine\\

L'artiste ne possède que l'art d'imiter \\
L'oeuvre d'art du pdv ontologique se situe à 2 degrés du réel\\

paradoxalement : comment un truc à peine vrai ait un tel pouvoir ? \\
L'art agit sur l'âme de manière irréfléchie \\
- Interdire la présence d'artistes dans la cité \\
- demander aux artistes de représenter des gens bons\\
- demander aux artistes de justifier leur art \\
\(\to\) romantisme, les gens produisent art + théorie de l'art 

Pour Platon, si on produit de l'art du pdv humain on ne développe pas 
une vision plus proche de la vérité\\

\noindent
L'Oeuvre : 1886\\

\noindent
XIXe siècle : \\
Romantisme vs Réalisme \\
(Delacroix) vs (Courbet)\\
pdv du sujet vs réel \\

L'art essaie de se renouveler, qq propositions : \\
- L'académisme (l'École des Beaux Arts ) \\
- L'impressionisme \\
1850 : peintres peignaient dans un atelier \\
Barbison : disait que les peintres devaient aller peindre dehors\\
pq aller voir ? pr que l'oeuvre d'art sort de l'illusoire vers 
la réalité.\\
Zola naturalisme : réalisme ++ \\
influence du milieu sur les individus + héritance \\
Sandoz == Zola \\
Les artistes sont représentés comme des êtres marginaux, comment percer ? \\
Début de la spéculation autour de l'art\\
Début de la presse, de la publicité \\
L'art : donne accès au vrai ce que l'on ne voit pas \\
On a tendance à conceptualiser quand on représente le monde mais pas l'artiste\\
L'art doté d'une telle puissance dans une oeuvre écrase\\

\noindent
Woolf : \\

Déterminer les conditions matérielles et sociales pour la création : \\
3 premiers chap : comprendres les conséquences matérielles et immatérielles favorables
à la création\\
La société n'est pas organisée pour favoriser l'émancipation des femmes\\
Il faut un monde où on ne passe pas son temps à se faire brimer\\
Le génie est une réunion de conditions socio-économiques \\
Modèles idéal de l'artiste : androgyne, qui voit plus le réel 
que nous \\

\noindent
Archi : Dubuche\\
Peintre : bcp\\
Musique : Wagner\\

\noindent
Peut on penser les arts ? \\
Utilisent tous un matériau, tous organisent ce matériau selon un principe
intellectuel et sensible.\\

\noindent
Peut on les classer ? \\
Kant : classer selon la notion espace et temps
Arts plastiques, archi \(\to\) se déploient dans espace\\
Musique, poésie... \(\to\) arts du temps \\

\noindent
Alain : L'art coupé en 2 :\\
- Dimension sociale/religieuse 
- Notion de collectif, danse, lectures théâtrales

artistes glorieux vs artistes ratées

Dissert : Éluard \\
Que signfie être inspiré ? 
- nature divine, modèle quelconque etc etc
L'artiste est capable d'inspirer à son tour \\
- Quels types d'inspiration existe il ? \\

\noindent
Faire schéma pour représenter les relations entre les concepts\\
On peut penser que l'artiste a créé pour susciter une réponse 
du public, même si il ne déclenche pas la réaction escomptée :\\
Plein Air, L'Enfant Mort...
L'artiste doit être intègre envers lui même avant tout 

\noindent
Avant de rédiger : \\
est ce que ma notion principales est présente pendant toute ma copie ? \\
est ce que ma copie articule et démontres bien toutes les notions et propositions de l'auteur du sujet dès le I,
pointe des problèmes au II, et y répond au III?\\
est ce qu'on a fait évaluer les relations entre les concepts sur ma copie ? \\

\noindent
première incidence : les notions doivent être sans cesse reprises \\
deuxième incidence : 

DS : intro, 1ere partie + plan détaillé de la pt2 et pt3

\noindent
Quels sont les critères du succès ? \\
Artistes glorieux : Ion, Flagerolles, Homère, Delacroix, Courbet, Shakespeare\\
Artistes pas glorieux : Claude, la soeur de Shakespeare, Bongrand\\

\noindent
Chez Zola : mauvais artiste est celui qui cherche à plaire, celui qui recycle les principes\\
Claude est un bon artiste : c'est juste qu'il va trop loin\\
Woolf : mauvais artiste est celui qui n'arrive pas à créer une oeuvre\\

\noindent \colorbox{green}{Exercice de dissertation} : \\

\noindent
Quand les artistes regardent une chose ils la voient pour elle et 
non plus pour eux 
ils ne perçoivent plus simplement en vue d'agir, 
ils perçoivent pour percevoir - pour rien, pour le plaisir
par un certain côté d'eux-même, soit par la conscience, 
soit par un de leurs sens, ils naissent 
\underline{détachés}\\
Bergson\\

\noindent Différence entre les 2\\ 

\noindent
\colorbox{green}{\underline{artiste}} = 
\colorbox{green}{\underline{perception}} , 
\colorbox{green}{PLAISIR}\(\rightarrow\) \fbox{innée} 
\(\rightarrow\) forme de la conscience, sensibilité\\
\(\rightarrow\) pas intéressée, \underline{utilitaire}\\
\(\rightarrow\) la perception n'est pas orientée vers l'acte\\
\(\rightarrow\) mais vers la contemplation 
\(\rightarrow\) voir la chose pour elle-même\\
\colorbox{green}{\underline{le non artiste}} = 
\colorbox{green}{perception} \\
\(\Rightarrow\) intéressée, \underline{utilitaire}, NAÏVE\\

\noindent 
les artistes \(\rightarrow\) \colorbox{green}{détaché} \begin{itemize}
  \item des autres \\
  \item de l'utilité \\
\end{itemize}
l'artiste est singularisé :\\
- par sa perception du monde \\
- est non utilitaire / (le + d'aspects possibles des choses)\\
\underline{conséquemment} \(\rightarrow\) cette manière de 
percevoir est à la base de la création artistique\\
La finalité de la perception n'est pas aussi pur \(\rightarrow\) 
pour réaliser une oeuvre d'art\\
"Un artiste c'est quelqu'un qui voit davantage" - Woolf\\
Un artiste c'est qqn qui a une vision, l'art naît dans la capacité
à voir et percevoir

\noindent
la perception différente de l'artiste le conduit à être en marge\\
 
\noindent la perception ne suffit pas à la création d'une oeuvre\\
après la création, est ce que l'artiste est toujours détaché ? \\
\(\rightarrow\) non, car en réalisant une oeuvre, il recrée un
lien avec nôtre monde \\
L'artiste ne né pas comme ça, il se construit de cette manière\\

\noindent
Chez Woolf\\
Un artiste n'est pas en dehors de tout cadre social\\
Mais il y a des points que cela ne peut expliquer \\
Cadres social condition nécessaires mais pas suffisante\\
L'artiste est celui : 
\begin{itemize}
  \item qui arrive à voir le réel plus que les autres\\
  \item plus en lien avec le réel, non artistes mettent une distance entre eux et le réel \\ 
  \item il parvient à mieux saisir son rapport au réel 
\end{itemize}

Part de mystère : une part des artistes qui leur est propre, que Woolf nous en dit pas plus 

Si dans l'Oeuvre de Zola, Chaine est un artiste qui n'a aucun talent. 
Alors cela signifie que les artistes ont un talent. \\
Style : signe d'un bon artiste \\

Pour Zola : le coeur du mystère de la création artistique c'est le style\\
Vrais artistes : Sandoz, Claude (Courbet, Delacroix)\\
Des gens qui ont une philosophie de l'Homme\\

Pour Platon : la questiond de Platon est la même que les deux autres\\
Pas de distinction entre l'artisanat et l'art en Grèce antique (l'art de la tekne) \\
Pour Platon ce n'est pas le cas\\
L'art n'est ni équivalent aux autres tekne, ni à un savoir théorique \\
Mais alors c'est quoi ? Ni theoria ni praxis\\ 

Ce à quoi on a habitude détermine ce qu'on voie : le public est 
"formatté" à apprécier un certain type d'art et il est réticent à 
découvrir quelque chose de nouveaux\\

Une des missions de l'art est de montrer les codes de la société
et il peut les critiquant en choquant\\

\noindent L'artiste révèle ce qui est et invente ce qui pourrait être - M. Gonod \\
Woolf pense que l'art peut révèler \\
L'art ne vise pas seulement une appréciation esthétique \\

\noindent
la création implique un degré de nouveauté\\
- La part du nouveau \\
- Les arts interprétatifs sont des arts \\ 

Différents modèles de création : \\
\noindent
\(1^{ère}\) dimension : création divine \\
C'est un modèle, dont on en déduit un modèle de création artistique\\

\noindent
\(2^{ème}\) création humaine\\
création qui traduit la logique des Hommes.\\
création \(\neq\) fabrication \\

Part de surprise dans la création, on ne peux pas réaliser un concept 
à perfection \\

\noindent 2 types d'imaginations : \\
- reproductrice \(\Leftrightarrow\) \\
- créatrice : 







\end{document}
